% !TeX root = ../main.tex

\chapter{Proposed Methods}
\label{chap:methods}

\section{Overall Pipeline}
\label{sec:pipeline}

Figure-conceptually, the proposed system is a four-stage pipeline. A raw hand-rotation video is first passed through MediaPipe to obtain a temporal sequence of 21 hand-keypoint coordinates per hand. From the resulting skeletal time series, three families of kinematic descriptors are extracted: \emph{frequency}, \emph{clarity} (rhythmic regularity), and \emph{amplitude}. The resulting high-dimensional feature vector is then reduced by one of three feature-selection strategies, namely ANOVA F-value, L1-regularized Logistic Regression, and XGBoost gain-based importance. Finally, the selected features are fed into one of four machine-learning models: Logistic Regression, Random Forest, XGBoost, or an Adaptive Graph Convolutional Network (AGCN). The output is a single binary label indicating whether the subject is PD or healthy.

\begin{table}[H]
    \centering
    \caption{Stages of the proposed pipeline}
    \label{tab:pipeline_overview}
    \small
    \setlength{\tabcolsep}{8pt}
    \begin{tabularx}{\linewidth}{>{\raggedright\arraybackslash}p{0.28\linewidth} >{\raggedright\arraybackslash}X}
        \hline
        Stage & Module / Output \\
        \hline
        Skeleton extraction   & MediaPipe Hands $\rightarrow$ 21-keypoint 3D skeleton per frame \\
        Feature extraction    & Frequency, Clarity, Amplitude \\
        Feature selection     & ANOVA F-value / L1-Logistic Regression / XGBoost Importance \\
        Machine learning      & Logistic Regression / Random Forest / XGBoost / AGCN \\
        Output                & PD-versus-healthy binary classification \\
        \hline
    \end{tabularx}
\end{table}

The remainder of this chapter describes each stage in detail, in five parts for feature extraction (Sections~\ref{sec:fe1}--\ref{sec:fe5}), followed by feature selection (Section~\ref{sec:fs}) and the joint-correlation / AGCN module (Section~\ref{sec:joint_corr}).

\section{Feature Extraction~(1/5): Frequency and Clarity}
\label{sec:fe1}

\subsection*{Observation}

Two qualitative observations motivate the frequency and clarity descriptors. First, severe PD patients tend to perform the pronation-supination task at a markedly slower pace; thus rotation \emph{frequency} carries a strong discriminative signal. Second, PD patients often exhibit irregular rhythm, with rotations that alternate between fast and slow within a single recording; thus the \emph{stability} of the rotation, captured below as a clarity index, is also informative.

\subsection*{Method}

\textbf{Polynomial detrend.} For each joint--axis trajectory, a low-frequency drift term is first fit by a polynomial and subtracted from the raw signal. This removes slow drifts introduced by gradual changes in hand position or camera framing. Conceptually,
\begin{equation}
    \text{Wave} + \text{Trend} \xrightarrow{\text{polynomial detrend}} \text{Detrend Result}.
\end{equation}

\textbf{Autocorrelation.} Let $r(\tau)$ denote the autocorrelation of the detrended signal. Two scalar descriptors are computed:
\begin{itemize}
    \item \textbf{Frequency.} The frame index of the first positive-lag peak of $r(\tau)$, which corresponds to the dominant period of rotation.
    \item \textbf{Clarity.} The height of the first peak divided by the value of $r(\tau)$ at lag zero. A clarity value close to one indicates a highly periodic signal, whereas a low clarity value reflects irregular rhythm.
\end{itemize}

\section{Feature Extraction~(2/5): Cycle Amplitude and Rolling Peak-to-Peak}
\label{sec:fe2}

\subsection*{Observation}

PD patients typically show a gradual reduction of movement amplitude over the course of a single recording, and instantaneous amplitudes may fluctuate from cycle to cycle. The amplitude features described in Sections~\ref{sec:fe2}--\ref{sec:fe3} are designed to capture these patterns through four complementary mathematical constructions.

\subsection*{Method}

\textbf{Cycle by cycle.} Iterating over every frame of the detrended trajectory, the algorithm flags each frame as a local maximum or minimum. Adjacent peaks and valleys are then paired, and the absolute amplitude of each rotation cycle is recorded, producing a derived time series whose value at each pair location is the instantaneous one-cycle amplitude.

\textbf{Rolling peak-to-peak.} A sliding window traverses the trajectory and, at each window position, records the difference between the local maximum and the local minimum within the window. The resulting time series captures the local range of motion without requiring exact peak/valley detection.

\section{Feature Extraction~(3/5): Rolling RMS and Hilbert Envelope}
\label{sec:fe3}

\textbf{Rolling RMS.} A sliding window computes the root-mean-square value of the (zero-mean) detrended signal within the window, providing a smoothed estimate of local movement energy. Among the four amplitude representations, this is the most robust to short-lived noise spikes.

\textbf{Hilbert envelope.} The instantaneous amplitude envelope is obtained from the analytic signal. The raw waveform is treated as the cosine component of an unknown amplitude--phase representation; the Hilbert transform produces the corresponding $90^{\circ}$-shifted (sine) component; and the magnitude of the resulting complex analytic signal yields the envelope. The envelope thus represents the instantaneous amplitude of the original oscillation, sharing the same phase progression but stripped of its oscillatory content.

\section{Feature Extraction~(4/5): Comparison of the Four Amplitude Representations}
\label{sec:fe4}

For each amplitude representation, three summary statistics are recorded: \emph{mean}, \emph{median}, and \emph{standard deviation}. Together, the four amplitude families therefore contribute $4 \times 3 = 12$ descriptors per joint axis.

The four representations are designed to be complementary. Table~\ref{tab:amplitude_comparison} summarizes the qualitative trade-offs that motivate their joint inclusion.

\begin{table}[H]
    \centering
    \caption{Comparison of the four amplitude representations}
    \label{tab:amplitude_comparison}
    \small
    \setlength{\tabcolsep}{4pt}
    \begin{tabularx}{\linewidth}{>{\raggedright\arraybackslash}p{0.22\linewidth} >{\raggedright\arraybackslash}X >{\raggedright\arraybackslash}X >{\raggedright\arraybackslash}X >{\raggedright\arraybackslash}X}
        \hline
        Property & Cycle by cycle & Rolling peak-to-peak & Rolling RMS & Hilbert envelope \\
        \hline
        Physical meaning   & Per-cycle displacement & Local range of motion & Local movement energy & Instantaneous envelope \\
        Noise sensitivity  & Highest                & High                  & Low                   & High \\
        Smoothness         & Low                    & Medium                & Highest               & High \\
        Movement aspect    & Single-cycle behaviour & Local activity        & Overall decay trend   & Overall decay trend \\
        \hline
    \end{tabularx}
\end{table}

\section{Feature Extraction~(5/5): Final Feature-Vector Dimensionality}
\label{sec:fe5}

Feature extraction is applied independently to every combination of hand, joint, and spatial axis. Per subject, the resulting feature vector therefore has the following dimensionality:
\begin{enumerate}
    \item Two hands, each with 21 joints and 3 spatial axes, yielding $2 \times 21 \times 3 = 126$ trajectories.
    \item Per trajectory, two autocorrelation features (frequency and clarity) and twelve amplitude features.
    \item Total dimensionality:
    \begin{equation}
        2 \times 21 \times 3 \times (2 + 4 \times 3) = 1764.
        \label{eq:feature_dim}
    \end{equation}
\end{enumerate}

\section{Feature Selection}
\label{sec:fs}

\subsection*{Motivation}

The feature dimensionality of $1{,}764$ greatly exceeds the number of available samples in either clinical cohort. Training a classifier directly on the full feature vector therefore tends to overfit, and the resulting model is hard to interpret. To mitigate this, three feature-selection methods are evaluated.

\subsection*{Methods}

\textbf{ANOVA F-value.} For each feature, the ratio between the between-class variance (PD vs healthy) and the within-class variance is computed. Features with the largest F-values are retained.

\textbf{L1-regularized Logistic Regression.} An L1 penalty is added to the Logistic Regression loss, which shrinks the coefficients of unimportant features all the way to zero. Features whose absolute coefficient is non-zero (or, equivalently, whose absolute coefficient is among the largest) are retained.

\textbf{XGBoost feature importance.} An XGBoost model is fit on the training fold and the gain-based importance of each feature, averaged over all decision-tree split points, is used as the selection criterion.

\subsection*{Qualitative Comparison}

The three methods differ in how they account for feature interactions and in how interpretable the resulting selection is. Table~\ref{tab:fs_methods_comparison} summarizes this contrast.

\begin{table}[H]
    \centering
    \caption{Qualitative comparison of the three feature-selection methods}
    \label{tab:fs_methods_comparison}
    \small
    \setlength{\tabcolsep}{6pt}
    \begin{tabularx}{\linewidth}{>{\raggedright\arraybackslash}p{0.32\linewidth} >{\raggedright\arraybackslash}X >{\raggedright\arraybackslash}X >{\raggedright\arraybackslash}X}
        \hline
        Method & ANOVA F-value & L1-Logistic Regression & XGBoost \\
        \hline
        Feature-interaction modelling & None              & Linear relationships    & Decision-tree combinations \\
        Interpretability              & Direct statistical significance & Clear linear weights    & Complex tree-based logic \\
        \hline
    \end{tabularx}
\end{table}

\section{Joint-Correlation Information: Pairwise Correlation and AGCN}
\label{sec:joint_corr}

\subsection*{Observation}

Some pairs of joints are not directly connected in the anatomical skeleton, but their motion may nevertheless be tightly coupled (or, conversely, abnormally decoupled) under PD. Encoding such inter-joint relationships explicitly can therefore reveal disease-related patterns that single-joint descriptors miss.

\subsection*{Pairwise Correlation Features}
\label{sec:joint_corr_corr}

A first, hand-crafted, way of encoding inter-joint relationships is to add the Pearson correlation between feature trajectories of pairs of joints as additional features. For each pair of joint axes $(i,j)$ with $i \neq j$, the correlation
\begin{equation}
    r_{ij} = \frac{\sum_{t} (s_i(t) - \bar s_i)(s_j(t) - \bar s_j)}
                  {\sqrt{\sum_{t} (s_i(t) - \bar s_i)^{2}} \, \sqrt{\sum_{t} (s_j(t) - \bar s_j)^{2}}}
\end{equation}
is computed over the entire recording. Within-hand and between-hand correlations are both retained. This descriptor, however, only carries information about \emph{two} joints at a time.

\subsection*{AGCN}
\label{sec:joint_corr_agcn}

A second, learned, way of encoding inter-joint relationships is to use a graph convolutional network. The AGCN reviewed in Section~\ref{sec:rw_agcn} trains an adaptive adjacency that captures motion-based joint relationships across \emph{all} joints in both hands jointly. Compared with the pairwise correlation descriptors above, the AGCN can in principle express higher-order, jointly-defined patterns of inter-joint coupling.

The two approaches are therefore presented as a deliberate contrast: pairwise correlation provides a simple hand-crafted baseline that only carries two-joint interactions, while AGCN learns a global, data-driven, multi-joint relational structure. Their respective contributions are evaluated empirically in Experiments~4 and~5 of Chapter~\ref{chap:results}.
